\section{Biological Plausibility：小脑能否实现 SDM 式读写}

理论对应之后，第一讲把问题推进到实现层：如果 attention-like retrieval 只需要高维激活、独立写入通道与 pooled readout，大脑是否已有类似电路？课程选择 cerebellum，不是声称小脑“就是 Transformer”，而是逐项寻找 address activation、value writing、superposition 与 majority-like pooling 的候选组件。

\subsection{抽象电路：地址通道与写入通道分开}

Circuitry 页从 abstract neuron 开始。输入 pattern 决定哪些地址 neuron firing；另一路 value signal 决定这些 neuron 应存什么；读取时 query 再次激活 neuron，各自输出叠加 value，最后由 pooled unit 汇总。这个 hetero-associative 结构与 attention 的 key/value 分离直接对应。

\lecturefigure{slide-36-circuitry-sdm.jpg}{实现 SDM 的抽象神经电路}{官方录像恢复幻灯片 V036；00:26:38}

\lecturefigure{slide-37-circuitry-mapping-final.jpg}{抽象电路映射到 cerebellar components 的最终状态}{官方录像恢复幻灯片 V037；00:28:22}

\begin{knowledgebox}{术语消化：cerebellar circuit 的候选角色}
\term{mossy fibers} 携带输入；\term{granule cells} 形成高维稀疏扩展；\term{parallel fibers} 把 granule activity 广泛送到 Purkinje cells；\term{climbing fibers} 提供强而相对独立的教学/写入信号；\term{Purkinje cells} 汇总大量输入并影响 deep cerebellar nuclei。课程是在功能层对齐这些角色。
\end{knowledgebox}

\subsection{小脑解剖与 SDM 操作逐项映射}

解剖页展示小脑重复而相对均一的微电路；最终 mapping 页把 abstract read/write 线路叠到具体 cell types。Granule cells 数量巨大且连接稀疏，可作为 hard locations；mossy input 决定激活；climbing fiber 的强信号可改变特定 Purkinje synapse；Purkinje cell 接收极大量 parallel-fiber input，提供 pooled output。

\lecturefigure{slide-38-cerebellum-anatomy.jpg}{Cerebellum 的神经元类型与连接结构}{官方录像恢复幻灯片 V038；00:31:04}

\lecturefigure{slide-39-sdm-cerebellum-mapping.jpg}{SDM read/write 与 cerebellar circuitry 的完整对照}{官方录像恢复幻灯片 V039；00:34:00}

\teachervoice{讲者说 70\% 左右的脑神经元位于 cerebellum，并提醒它可能不仅参与精细运动；但他同时把高阶认知证据描述为活跃研究。讲义保留这个探索性语气，不把 fMRI “亮起”改写成某种已确定的语言 attention 模块。}

\subsection{怎样从“可实现”走向“被证实”}

课程问答给出最重要的科学边界：即使电路可以实现 SDM，也不代表大脑实际使用它。真正检验需要在 associative task 中实时记录稀疏 neuron activation、synaptic update 与 behavior change，最好还能写入新关联并观察旧关联消退。讲者提出高速 calcium/voltage imaging 作为候选路线；这是一份实验设计，而不是已有结论。

\begin{warningbox}{三种不能跳过的证据差距}
第一，解剖连接允许某计算，不代表动力学会执行它；第二，局部读写相似，不代表全脑使用 Transformer training objective；第三，microzone 与 multi-head 的类比很诱人，但讲者明确说不应把类比推得太远。
\end{warningbox}

Q\&A 还区分 fast neural activity 与 long-term synaptic plasticity，并讨论 radius 的容量--fidelity 权衡：激活半径大，单次读出聚合更多 neuron、估计更平滑，却增加 pattern 干扰；半径小，区分度高，但噪声 query 可能激活太少位置。attention 中对应参数不是一个固定物理半径，而是 key/query norms 与 beta 共同产生的 effective sharpness。

\lecturefigure{slide-40-first-talk-thank-you.jpg}{第一讲总结、参考资料与致谢}{官方录像恢复幻灯片 V040；00:35:54}

\begin{importantbox}{第一讲的最强与最弱结论}
最强结论是：在明确条件下，SDM intersection read 与 softmax attention 存在可推导、可数值验证的近似。较弱但有启发的结论是：cerebellar circuit 拥有实现类似 address/value/readout 操作的组件。二者必须分开写，前者是数学证据，后者是神经科学假设。
\end{importantbox}

\subsection{本章小结}

小脑 mapping 展示了 attention-like associative memory 的候选生物实现：稀疏扩展、独立教学信号、分布式存储与 pooled readout 都能找到功能对应。但课堂最重要的 teacher voice 是“可实现不等于被验证”，实验需要直接追踪学习中的连接变化和行为结果。

